\setcounter{section}{4}


\inputaufgabe
{}
{





Es seien
\mathkor {} {{ \mathcal  F   }} {und} {{ \mathcal  G   }} {}
\definitionsverweis {Garben}{}{}
auf dem
\definitionsverweis {topologischen Raum}{}{}
$X$. Zeige, dass durch
\mathl{U \mapsto { \mathcal  F   }  \left(   U   \right) \times { \mathcal  G   }  \left(   U   \right)}{} mit den natürlichen
\definitionsverweis {Produktabbildungen}{}{}
eine Garbe auf $X$ gegeben ist.





}
{} {}

\inputaufgabe
{}
{





Es sei ${ \mathcal  G   }$ eine
\definitionsverweis {Garbe}{}{}
auf einem nicht
\definitionsverweis {zusammenhängenden Raum}{}{}
$X$ mit einer Zerlegung

\mavergleichskette
{\vergleichskette
{X
}
{ = }{ U \uplus V
}
{  }{
}
{    }{
}
{   }{
}
}
{}{}{}
in
\definitionsverweis {disjunkte}{}{}
nichtleere offene Mengen. Zeige

\mavergleichskette
{\vergleichskette
{ { \mathcal  G   }  \left(   X    \right)
}
{ = }{ { \mathcal  G   }  \left(   U    \right) \times { \mathcal  G   }  \left(   V    \right)
}
{  }{
}
{    }{
}
{   }{
}
}
{}{}{.}





}
{} {}

\inputaufgabe
{}
{





Es sei $X$ ein
\definitionsverweis {topologischer Raum}{}{}
mit einer Zerlegung

\mavergleichskette
{\vergleichskette
{X
}
{ = }{ Y \uplus Z
}
{  }{
}
{    }{
}
{   }{
}
}
{}{}{}
in disjunkte offene nichtleere Teilmengen. Es sei ${ \mathcal  G   }$ eine
\definitionsverweis {Garbe}{}{}
auf $Y$ und ${ \mathcal  H   }$ eine Garbe auf $Z$. Zeige, dass durch

\mavergleichskettedisp
{\vergleichskette
{ { \mathcal  F   }  \left(   U   \right)
}
{ =} { { \mathcal  G   }  \left(   U \cap Y   \right)  \times { \mathcal  H   }  \left(   U \cap Z   \right)
}
{ } {
}
{ } {
}
{ } {
}
}
{}{}{}
für

\mavergleichskette
{\vergleichskette
{U
}
{ \subseteq }{ X
}
{  }{
}
{    }{
}
{   }{
}
}
{}{}{}
eine Garbe auf $X$ definiert ist.





}
{} {}

\inputaufgabe
{}
{





Es sei $X$ ein
\definitionsverweis {Hausdorffraum}{}{}
mit zumindest zwei Punkten und es sei

\mavergleichskette
{\vergleichskette
{M
}
{ \neq }{ \emptyset
}
{  }{
}
{    }{
}
{   }{
}
}
{}{}{.}
eine Menge. Zeige, dass die
\definitionsverweis {konstante Prägarbe}{}{}
zu $M$ keine
\definitionsverweis {Garbe}{}{}
ist.





}
{} {}

\inputaufgabe
{}
{





Zeige, dass die Einschränkung einer
\definitionsverweis {Garbe}{}{}
auf eine
\definitionsverweis {offene Teilmenge}{}{}

\mavergleichskette
{\vergleichskette
{U
}
{ \subseteq }{ X
}
{  }{
}
{    }{
}
{   }{
}
}
{}{}{}
eine Garbe ist.





}
{} {}

\inputaufgabe
{}
{





Zeige, dass der
\definitionsverweis {Halm}{}{}
der Garbe der
\definitionsverweis {holomorphen Funktionen}{}{}
in

\mavergleichskette
{\vergleichskette
{ 0
}
{ \in }{ {\mathbb C}
}
{  }{
}
{    }{
}
{   }{
}
}
{}{}{}
gleich dem Ring der konvergenten Potenzreihen in einer Variablen ist.





}
{} {}

\inputaufgabe
{}
{





Es sei
\maabb {\varphi} { { \mathcal  F   } } { { \mathcal  G   }
} {}
ein
\definitionsverweis {Garbenmorphismus}{}{}
zwischen
\definitionsverweis {Garben}{}{}
auf einem
\definitionsverweis {topologischen Raum}{}{}
$X$ und es sei
\maabb {\varphi_U} { { \mathcal  F   }  \left(   U   \right) } { { \mathcal  G   }  \left(   U   \right)
} {}
\definitionsverweis {surjektiv}{}{}
für jede offene Menge

\mavergleichskette
{\vergleichskette
{U
}
{ \subseteq }{ X
}
{  }{
}
{    }{
}
{   }{
}
}
{}{}{.}
Zeige, dass dann auch jede
\definitionsverweis {Halmabbildung}{}{}
\maabb {\varphi_P} { { \mathcal  F   }_P } { { \mathcal  G   }_P
} {}
surjektiv ist.





}
{} {}

\inputaufgabe
{}
{





Es sei ${ \mathcal  G   }$ eine
\definitionsverweis {Garbe von kommutativen Gruppen}{}{}
auf einem
\definitionsverweis {topologischen Raum}{}{}
$X$. Zeige

\mavergleichskette
{\vergleichskette
{ { \mathcal  G   }  \left(   \emptyset   \right)
}
{ = }{ 0
}
{  }{
}
{    }{
}
{   }{
}
}
{}{}{,}
also die triviale Gruppe ist.





}
{} {}

\inputaufgabe
{}
{





Es sei $X$ ein
\definitionsverweis {topologischer Raum}{}{,}

\mavergleichskette
{\vergleichskette
{ P
}
{ \in }{ X
}
{  }{
}
{    }{
}
{   }{
}
}
{}{}{}
ein Punkt und $G$ eine
\definitionsverweis {kommutative Gruppe}{}{.}
Wir betrachten die Zuordnung

\mathdisp {U \longmapsto { \mathcal  G   }  \left(   U   \right) \defeq  \begin{cases} G  \, , \text{ falls } P \in U \\   0 \text{ sonst} \, .  \end{cases}} {  }

mit den naheliegenden Restriktionsabbildungen zu offenen Teilmengen

\mavergleichskette
{\vergleichskette
{ V
}
{ \subseteq }{ U
}
{  }{
}
{    }{
}
{   }{
}
}
{}{}{.}
\aufzaehlungdrei{Zeige, dass ${ \mathcal  G   }  \left(   U   \right)$ eine
\definitionsverweis {Garbe von kommutativen Gruppen}{}{}
ist.
}{Bestimme den
\definitionsverweis {Halm}{}{}
von ${ \mathcal  G   }$ im Punkt $P$.
}{Es sei nun $P$ ein abgeschlossener Punkt. Besitmme die Halm für jeden Punkt

\mavergleichskette
{\vergleichskette
{ Q
}
{ \neq }{ P
}
{  }{
}
{    }{
}
{   }{
}
}
{}{}{.}
}





}
{} {}

Die in der vorstehenden Aufgabe konstruierte Garbe nennt man \stichwort {Wolkenkratzergarbe} {}
\zusatzklammer {zu $G$ im Punkt $P$} {} {.}
